\chapter{Interface Specification}
\label{ch:interface}

This chapter specifies every surface, every region and every control. Each control carries a
stable identifier used by the implementation, the automated interface tests and the
accessibility layer. A control that is not in this chapter does not ship.

\section{Design system}

\subsection{Tokens}

\begin{table}[h]
\centering
\small
\begin{tabularx}{\textwidth}{l l X}
\toprule
\textbf{Token} & \textbf{Value} & \textbf{Applied to} \\
\midrule
space.1 to space.8 & 2, 4, 6, 8, 12, 16, 24, 32 & All padding and gaps, no arbitrary values \\
radius.sm / md / lg & 4, 6, 10 & Controls, cards, panels \\
control.height.sm & 28 & Compact toolbar controls \\
control.height.md & 34 & Default buttons, inputs \\
control.height.lg & 42 & Primary actions, empty state actions \\
text.xs to text.xl & 11, 12, 13, 15, 19 & Labels, body, headings \\
icon.sm / md / lg & 14, 16, 20 & Inline, control, feature \\
elevation.0 to 3 & Flat, subtle, raised, floating & Panels, popovers, dialogs, drag \\
\bottomrule
\end{tabularx}
\caption{Design tokens. Sizes are in logical pixels at a scale factor of one.}
\end{table}

\subsection{Rules that are not negotiable}

\begin{itemize}
  \item Every control carries a visible text label unless it is in a dense toolbar, and even
  then it carries an accessible name and a tooltip. Icon only interfaces are a discoverability
  failure.
  \item Minimum interactive target is thirty four logical pixels in the primary flow.
  \item Sentence case everywhere. No capitalised words in interface copy.
  \item Every destructive action is either reversible or confirmed, never both silent and
  permanent.
  \item Every long running action shows determinate progress when a total is known, and an
  elapsed counter when it is not.
  \item Colour never carries meaning alone. Severity is colour plus a glyph plus a word.
\end{itemize}

\section{Window anatomy}

\begin{figure}[h]
\centering
\begin{tikzpicture}[
  reg/.style={draw=xyraline,fill=xyrapanel,font=\scriptsize,align=center},
  hi/.style={draw=xyralime,line width=0.9pt,fill=white,font=\scriptsize,align=center}
]
\draw[draw=xyragrey,line width=1pt,rounded corners=3pt] (0,0) rectangle (14,9.4);
\node[reg,minimum width=13.6cm,minimum height=0.7cm] at (7,9.0) {Title bar: project, branch, mission pill, model pill, window controls};
\node[reg,minimum width=1.0cm,minimum height=6.9cm] at (0.55,5.0) {A\\c\\t\\i\\v\\i\\t\\y};
\node[reg,minimum width=2.6cm,minimum height=6.9cm] at (2.45,5.0) {Left dock\\(Explorer,\\Search,\\Source control,\\Atlas)};
\node[hi,minimum width=5.6cm,minimum height=4.4cm] at (6.6,6.25) {Editor group\\tabs, breadcrumb, gutter,\\code, inline agent affordances};
\node[reg,minimum width=5.6cm,minimum height=2.3cm] at (6.6,2.45) {Bottom dock\\Terminal, Problems, Verification Deck, Output};
\node[hi,minimum width=3.6cm,minimum height=6.9cm] at (11.9,5.0) {Right dock\\Mission Control,\\Agent Theatre,\\Context Inspector,\\Chat};
\node[reg,minimum width=13.6cm,minimum height=0.55cm] at (7,1.05) {Status bar: panel toggles with labels, diagnostics, verification, tokens, cost, EKG};
\end{tikzpicture}
\caption{Window regions. The right dock is where autonomous work is observed and steered.}
\end{figure}

\section{Title bar}

\begin{longtable}{@{}p{0.30\textwidth} p{0.24\textwidth} p{0.36\textwidth}@{}}
\toprule
\textbf{Control} & \textbf{Label} & \textbf{Behaviour} \\
\midrule
\endhead
ctl.title.project & Project name & Opens the recent projects menu \\
ctl.title.branch & Branch name with dirty dot & Opens branch switcher with search \\
ctl.title.mission & Mission pill, for example \ctl{7/19} & Present only during a mission. Click focuses Mission Control. Shows a determinate ring \\
ctl.title.model & Active model and role routing & Opens the routing popover: implementer, reviewer, planner, embedding \\
ctl.title.share & Share & Opens the collaboration menu when enabled \\
ctl.title.zoom & Window controls & Platform native \\
\bottomrule
\caption{Title bar controls.}
\end{longtable}

\section{Activity bar}

Vertical, icon plus short label, ordered and reorderable. Each entry toggles a dock panel and
shows a badge when it has unseen state. Five entries open the primary surfaces named in
Chapter~\ref{ch:thesis}; the remainder open supporting panels that live inside the Workbench
surface rather than constituting surfaces of their own.

\begin{table}[h]
\centering
\small
\begin{tabularx}{\textwidth}{l l l X}
\toprule
\textbf{Control} & \textbf{Label} & \textbf{Badge} & \textbf{Opens} \\
\midrule
ctl.activity.explorer & Files & Modified count & Explorer \\
ctl.activity.search & Search & Match count & Search \\
ctl.activity.vcs & Git & Changed count & Source control \\
ctl.activity.mission & Mission & Progress ring & Mission Control \\
ctl.activity.agents & Agents & Active lanes & Agent Theatre \\
ctl.activity.context & Context & Budget percent & Context Inspector \\
ctl.activity.verify & Verify & Failing count & Verification Deck \\
ctl.activity.atlas & Atlas & None & Atlas \\
ctl.activity.ext & Extensions & Update count & Extensions \\
\bottomrule
\end{tabularx}
\caption{Activity bar entries.}
\end{table}

\section{Workbench: the editor surface}

\subsection{Tab strip}

\begin{longtable}{@{}p{0.30\textwidth} p{0.62\textwidth}@{}}
\toprule
\textbf{Control} & \textbf{Behaviour} \\
\midrule
\endhead
ctl.tab.item & File name, dirty dot, diagnostic underline, close on middle click \\
ctl.tab.agentmark & A small marker on tabs whose content was written by an agent and not yet reviewed by a human \\
ctl.tab.pin & Pin, keeps the tab through bulk close \\
ctl.tab.split & Split right, split down \\
ctl.tab.overflow & Overflow list with search when tabs exceed the width \\
ctl.tab.newfile & New file in this group \\
\bottomrule
\caption{Tab strip controls.}
\end{longtable}

\subsection{Gutter and inline affordances}

\begin{longtable}{@{}p{0.30\textwidth} p{0.62\textwidth}@{}}
\toprule
\textbf{Control} & \textbf{Behaviour} \\
\midrule
\endhead
ctl.gutter.line & Line numbers, relative mode optional \\
ctl.gutter.fold & Fold and unfold by syntax node \\
ctl.gutter.vcs & Added, modified and removed markers, click opens the inline diff \\
ctl.gutter.blame & Inline blame on hover, showing the commit and its message \\
ctl.gutter.agent & Marker on ranges an agent changed in the current mission, with the ticket identifier \\
ctl.inline.ask & Inline agent prompt at the cursor, opened with a shortcut \\
ctl.inline.explain & Explain this symbol, uses structural context rather than a raw dump \\
ctl.inline.impact & Show impact for the symbol under the cursor, opens the closure in Atlas \\
ctl.inline.verify & Verify this file now, runs mapped tests only \\
ctl.inline.accept & Accept the agent range, clears the marker \\
ctl.inline.revert & Revert the agent range to the pre mission state \\
\bottomrule
\caption{Gutter and inline controls.}
\end{longtable}

\subsection{Breadcrumb and minimap}

\begin{itemize}
  \item \ctl{ctl.breadcrumb.segment} navigates by structural path, module then type then function,
  each segment opening a filtered symbol list.
  \item \ctl{ctl.minimap.toggle} switches the minimap between off, code and heat, where heat colours
  the map by how much agent attention each region received in the current mission.
\end{itemize}

\section{Mission Control}

\begin{figure}[h]
\centering
\begin{tikzpicture}[
  reg/.style={draw=xyraline,fill=xyrapanel,font=\scriptsize,align=center},
  hi/.style={draw=xyralime,line width=0.9pt,fill=white,font=\scriptsize,align=center}
]
\draw[draw=xyragrey,line width=1pt,rounded corners=3pt] (0,0) rectangle (7.6,10.2);
\node[reg,minimum width=7.2cm,minimum height=1.1cm] at (3.8,9.5) {Objective, status, mission identifier};
\node[hi,minimum width=7.2cm,minimum height=1.0cm] at (3.8,8.2) {Progress bar, tickets done over total, elapsed};
\node[reg,minimum width=7.2cm,minimum height=0.85cm] at (3.8,7.05) {Start, Pause, Stop, Resume, Replan};
\node[reg,minimum width=7.2cm,minimum height=0.85cm] at (3.8,6.05) {Budget: sessions, hours, consecutive failures};
\node[hi,minimum width=7.2cm,minimum height=3.4cm] at (3.8,4.05) {Ticket list\\identifier, title, state glyph,\\attempts, commit, duration};
\node[reg,minimum width=7.2cm,minimum height=1.6cm] at (3.8,1.55) {Quarantine drawer\\failure text and the split proposal};
\node[reg,minimum width=7.2cm,minimum height=0.5cm] at (3.8,0.45) {Open journal, open cockpit};
\end{tikzpicture}
\caption{Mission Control layout.}
\end{figure}

\begin{longtable}{@{}p{0.28\textwidth} p{0.18\textwidth} p{0.46\textwidth}@{}}
\toprule
\textbf{Control} & \textbf{Label} & \textbf{Behaviour} \\
\midrule
\endhead
ctl.mission.objective & Objective field & Multi line, editable before start, read only during a run \\
ctl.mission.plan & Plan & Produces the ticket graph without executing, so it can be reviewed \\
ctl.mission.start & Start & Plans if needed, then runs to completion \\
ctl.mission.pause & Pause & Halts after the ticket in flight, state preserved \\
ctl.mission.resume & Resume & Continues from durable state, also available after a crash \\
ctl.mission.stop & Stop & Halts and marks the mission stopped \\
ctl.mission.background & Run in background & Detaches the supervisor so it survives closing the window \\
ctl.mission.budget & Budget & Popover with sessions, hours, attempts per ticket, consecutive failures \\
ctl.mission.ticket & Ticket row & Click opens the ticket detail: brief, diff, verification output, transcript \\
ctl.mission.retry & Retry ticket & Requeues a quarantined ticket, optionally with a different vendor \\
ctl.mission.split & Split ticket & Sends the ticket to the replanner and inserts the children \\
ctl.mission.skip & Skip ticket & Marks it abandoned and unblocks its dependants where possible \\
ctl.mission.edit & Edit ticket & Modifies title, scope and files before it runs \\
ctl.mission.reorder & Drag handle & Reorders within the constraints of the dependency graph \\
ctl.mission.addticket & Add ticket & Inserts a manual ticket into the graph \\
ctl.mission.journal & Open journal & Opens the decision journal filtered to this mission \\
ctl.mission.cockpit & Open cockpit & Opens the end of mission summary and the approval gate \\
\bottomrule
\caption{Mission Control controls.}
\end{longtable}

\subsection{Ticket states}

\begin{table}[h]
\centering
\small
\begin{tabularx}{\textwidth}{l l X}
\toprule
\textbf{State} & \textbf{Glyph} & \textbf{Meaning} \\
\midrule
pending & Hollow circle & Waiting for its dependencies \\
ready & Outlined circle & Dependencies satisfied, queued \\
running & Filled ring, animated & A session is executing it now \\
verifying & Half filled ring & Session finished, tests executing in a snapshot \\
done & Check & Verified and committed, commit hash shown \\
retry & Rotating arrow & Failed, queued again with a different vendor \\
split & Branch glyph & Replaced by children \\
quarantined & Warning triangle & Exhausted its policy, needs a human \\
skipped & Dash & Abandoned deliberately \\
\bottomrule
\end{tabularx}
\caption{Ticket state vocabulary. Every state has a glyph, a colour and a word.}
\end{table}

\section{Agent Theatre}

\begin{figure}[h]
\centering
\begin{tikzpicture}[
  lane/.style={draw=xyraline,fill=xyrapanel,minimum width=10.6cm,minimum height=0.85cm,font=\scriptsize,align=left},
  hi/.style={draw=xyralime,line width=0.9pt,fill=white,minimum width=10.6cm,minimum height=0.85cm,font=\scriptsize,align=left}
]
\draw[draw=xyragrey,line width=1pt,rounded corners=3pt] (0,0) rectangle (11.4,6.6);
\node[font=\scriptsize\bfseries] at (5.7,6.15) {Agent Theatre: live lanes, artefacts, cost};
\node[lane] at (5.7,5.4) {\hspace{2mm} Planner \quad idle \quad last: ticket graph, 19 nodes};
\node[hi] at (5.7,4.45) {\hspace{2mm} Implementer \quad running ticket 7 \quad 42s \quad 8.1k tokens};
\node[lane] at (5.7,3.5) {\hspace{2mm} Reviewer \quad queued \quad lenses: correctness, security, conventions};
\node[hi] at (5.7,2.55) {\hspace{2mm} Verifier \quad running \quad snapshot tests \quad 14/38 passed};
\node[lane] at (5.7,1.6) {\hspace{2mm} Repairer \quad idle \quad last: fixed a null guard in ticket 5};
\node[lane] at (5.7,0.65) {\hspace{2mm} Scribe \quad idle \quad last: journal entry for ticket 6};
\end{tikzpicture}
\caption{Agent Theatre. One lane per role from Table~\ref{tab:roles}, each showing its current
act and its last artefact. Architect and Replanner lanes appear only while those roles are
active, so an idle mission shows six lanes rather than eight.}
\end{figure}

\begin{longtable}{@{}p{0.28\textwidth} p{0.64\textwidth}@{}}
\toprule
\textbf{Control} & \textbf{Behaviour} \\
\midrule
\endhead
ctl.theatre.lane & Expands to the lane's transcript, tool calls and produced artefacts \\
ctl.theatre.artifact & Opens the artefact: a diff, a verdict, a screenshot, a report \\
ctl.theatre.kill & Terminates the session in that lane and marks its ticket failed \\
ctl.theatre.vendor & Switches the vendor for that role for subsequent sessions \\
ctl.theatre.follow & Follows the active lane automatically as work moves between roles \\
ctl.theatre.cost & Toggles the cost column between tokens, requests and elapsed \\
ctl.theatre.graph & Switches from lanes to a node graph showing which role handed work to which \\
ctl.theatre.filter & Filters lanes by role or by ticket \\
\bottomrule
\caption{Agent Theatre controls.}
\end{longtable}

\section{Context Inspector}

\begin{longtable}{@{}p{0.28\textwidth} p{0.64\textwidth}@{}}
\toprule
\textbf{Control} & \textbf{Behaviour} \\
\midrule
\endhead
ctl.context.budget & Budget bar: used versus available, segmented by tier, with the token count \\
ctl.context.item & One row per included item: path, tier, tokens, why it was included \\
ctl.context.why & Explains the inclusion reason: structural closure, pinned, similarity fallback, skill \\
ctl.context.pin & Pins an item so the solver must always include it \\
ctl.context.exclude & Excludes an item for this workspace and writes the rule to settings \\
ctl.context.expand & Raises an item's tier, for example from signatures to full body \\
ctl.context.collapse & Lowers an item's tier to reclaim budget \\
ctl.context.heat & Heat map over the file tree, coloured by accumulated agent attention \\
ctl.context.skills & Lists the skills injected into the current session with their source \\
ctl.context.integrations & Shows which expansion integrations contributed, with a capability toggle each \\
ctl.context.compression & Shows the compression ratio and the escalation rate for the session \\
ctl.context.export & Exports the exact assembled context for reproduction and debugging \\
\bottomrule
\caption{Context Inspector controls.}
\end{longtable}

\section{Verification Deck}

\begin{longtable}{@{}p{0.28\textwidth} p{0.64\textwidth}@{}}
\toprule
\textbf{Control} & \textbf{Behaviour} \\
\midrule
\endhead
ctl.verify.run & Verify now: snapshots the change set and runs the mapped tests \\
ctl.verify.full & Full suite: runs everything rather than the mapped subset \\
ctl.verify.result & A run: command, exit code, duration, failing tests, output tail \\
ctl.verify.rerun & Reruns a single failing test in isolation \\
ctl.verify.repair & Hands the failure to the repairer inside the snapshot \\
ctl.verify.visual & Renders the configured page and returns a visual verdict \\
ctl.verify.compare & Compares the render against a design export \\
ctl.verify.qa & Runs an adversarial sweep against the running application \\
ctl.verify.flags & Security and cost flags with severity, file and line \\
ctl.verify.policy & Opens the verification policy for this repository \\
ctl.verify.history & Previous runs with their verdicts, filterable by ticket \\
\bottomrule
\caption{Verification Deck controls.}
\end{longtable}

\section{Atlas}

\begin{longtable}{@{}p{0.28\textwidth} p{0.64\textwidth}@{}}
\toprule
\textbf{Control} & \textbf{Behaviour} \\
\midrule
\endhead
ctl.atlas.mode & Switches between topology, impact, ownership and history views \\
ctl.atlas.filter & Filters nodes by path, language or module \\
ctl.atlas.focus & Focuses a node and dims everything outside its neighbourhood \\
ctl.atlas.depth & Sets transitive depth for impact, from one to unbounded \\
ctl.atlas.open & Opens the file at the definition \\
ctl.atlas.mission & Creates a mission scoped to the current selection \\
ctl.atlas.fleet & Switches the scope from this repository to the registered fleet \\
ctl.atlas.export & Exports the current view as a vector image for documentation \\
\bottomrule
\caption{Atlas controls.}
\end{longtable}

\section{Status bar}

\begin{table}[h]
\centering
\small
\begin{tabularx}{\textwidth}{l X}
\toprule
\textbf{Control} & \textbf{Behaviour} \\
\midrule
ctl.status.panels & Labelled toggles: Files, Search, Git, Mission, Agents, Context, Verify, Terminal \\
ctl.status.diag & Errors and warnings, click opens Problems \\
ctl.status.verify & Last verification verdict with a glyph, click opens the Deck \\
ctl.status.tokens & Tokens used in the current session and the budget remaining \\
ctl.status.cost & Session cost when the provider reports it, otherwise request count \\
ctl.status.ekg & Activity pulse driven by real agent state: idle, thinking, verifying, failing \\
ctl.status.position & Cursor position, selection size, indentation, encoding, language \\
\bottomrule
\end{tabularx}
\caption{Status bar controls.}
\end{table}

\section{Command palette and decisions}

The palette is the universal entry point and is required to reach every command in this
document. It supports three modes: commands, symbols with a leading hash, and files with a
leading slash. Missions are startable directly from it with a leading exclamation mark.

Decision cards replace terminal questions. When an agent reaches a genuine fork it emits a
choice, and the interface renders two or three cards side by side, each carrying the option
name, its measured cost, its structural consequence and a code sketch. Selecting a card
records the decision in the journal with its rationale and resumes the mission.

\section{First run}

\begin{enumerate}
  \item Welcome: product statement, open a project, clone from a provider.
  \item Model routing: choose the implementer, the reviewer, the planner and the embedding
  model, with a local only option offered explicitly.
  \item Sign in to the chosen providers, each with its own state indicator.
  \item Index the project, showing determinate progress and what is being built.
  \item Optional fleet registration.
  \item A guided first mission on a scoped objective, so the developer sees the whole loop
  before trusting it with real work.
\end{enumerate}

\section{Keymap}

\begin{table}[h]
\centering
\small
\begin{tabularx}{\textwidth}{l X}
\toprule
\textbf{Binding} & \textbf{Command} \\
\midrule
mod+shift+p & Command palette \\
mod+p & Files \\
mod+shift+f & Search across the project \\
mod+i & Inline agent prompt at the cursor \\
mod+m & Toggle Mission Control \\
mod+shift+m & Start or resume the mission \\
mod+alt+m & Stop the mission \\
mod+g & Toggle Agent Theatre \\
mod+k then c & Toggle Context Inspector \\
mod+shift+v & Verify now \\
mod+alt+v & Full verification \\
mod+shift+a & Toggle Atlas \\
mod+alt+i & Impact for the symbol under the cursor \\
mod+j & Toggle terminal \\
mod+enter & Accept the agent range under the cursor \\
mod+backspace & Revert the agent range under the cursor \\
\bottomrule
\end{tabularx}
\caption{Default bindings. Every one is rebindable and every command is reachable without one.}
\end{table}

\section{States that must be designed, not improvised}

\begin{longtable}{@{}p{0.26\textwidth} p{0.66\textwidth}@{}}
\toprule
\textbf{State} & \textbf{Requirement} \\
\midrule
\endhead
Empty project & Explains what to do next with a primary action, never a blank pane \\
Indexing & Determinate progress, the product remains usable, agents wait rather than guess \\
No model configured & Blocks agent surfaces with a direct route to routing settings \\
Provider unauthenticated & Inline sign in with the exact command, never a silent failure \\
Mission running & Every surface shows it, and closing the window does not kill it \\
Verification failing & Failure is the loudest thing on screen, with the exact command to reproduce \\
Quarantined ticket & Failure text and the proposed split, with one click to accept the split \\
Offline & Local capabilities keep working and remote ones are visibly disabled \\
Large repository & Progressive indexing, structural queries available before embeddings finish \\
Conflicting edits & The developer's edit always wins, the agent range is marked stale and reverifies \\
\bottomrule
\caption{Designed states.}
\end{longtable}
